\section{Security Analysis}
\label{sec:security}

The evaluation leaves a question the accuracy numbers cannot answer: if the gateways and the aggregator are exposed to adversaries, how much of the measured benefit survives, and which parts of the system are actually defended? The answer depends on whom the system trusts, so the adversaries and assumptions come first. They are stated here once; Section~\ref{sec:system} describes the mechanisms without repeating them.

\subsection{Threat Model}
\label{sec:sec-threat}
The deployment is two or more edge gateways, each with a Cortex-A72-class CPU, no GPU and a standard-library Python runtime, and one cloud aggregator. Request payloads are private to the site and are never sent to the aggregator. Four adversaries are considered.

\begin{enumerate}
  \item \textbf{A1, application-layer attacker.} Sends injection and cross-site-scripting payloads through a gateway toward upstream services. The gateway classifier is the control under test.
  \item \textbf{A2, network attacker.} Observes, replays, reorders or forges messages between gateway and aggregator, and may record traffic indefinitely. A2 holds no enrolled secret.
  \item \textbf{A3, Sybil attacker.} Reaches the aggregator and tries to register or impersonate identities in order to outnumber honest peers.
  \item \textbf{A4, compromised gateway.} Holds a valid enrolled secret and uploads arbitrary parameters; it is a Byzantine client, and up to $f$ such nodes are assumed.
\end{enumerate}

The aggregator is trusted to run the aggregation rule and to keep enrolled secrets, but receives no raw payloads. Out of scope are a compromised aggregator and adaptive input-space evasion of the classifier (Section~\ref{sec:sec-limits}). Poisoning optimised against the aggregation rule~\cite{fang2020} is examined only in simulation (Section~\ref{sec:eval-byzantine}). The gateway does not authenticate the clients whose requests it inspects. A1 is handled by classification alone, which is why the paper treats the gateway as a payload-inspection enforcement point inside a zero-trust deployment and not as a zero-trust architecture.

The controls are ordered by what each makes possible for the next. Authentication bounds the number of identities and thereby makes $f$ finite. Admission checks drop malformed or grossly scaled updates before aggregation. Robust aggregation then contains the $f$ Byzantine peers that remain. Without authentication the last layer is void, because an attacker who mints identities controls the median regardless of its guarantee.

\subsection{Poisoning Defence}
\label{sec:sec-poison}
Coordinate-wise median and $\beta$-trimmed mean~\cite{yin2018} (see also Krum~\cite{blanchard2017}) resist Byzantine updates only while honest clients outnumber Byzantine ones in every coordinate:
\begin{equation}
  K \ge 2f + 1 .
  \label{eq:security:byz-bound}
\end{equation}
On the two-gateway testbed $K=2$ gives $f=0$, since the median of two values is their mean and a single Byzantine node is half the population. The aggregator therefore computes $f=\lfloor (K-1)/2\rfloor$ from the number of distinct enrolled peers that have reported and returns it with the model, so that no gateway is told it is protected when it is not. Updates are keyed by client identity and only the latest per peer is kept, so the median runs across peers. An earlier version used a window of the last ten uploads regardless of sender, which is a median over time; a persistent attacker then simply fills a fixed share of the window.

The $\ell_2$-norm gate rejects with HTTP~409 any update whose norm differs from the reference by more than a factor of $4$ (Section~\ref{sec:sys-agg}). It stops gross rescaling. It does not detect sign-flips of ordinary magnitude, nor an update shaped to stay just inside the ratio.

Section~\ref{sec:eval-byzantine} reports two simulations. The first, a bare-aggregator ablation with seven nodes and a $-12\times$ sign flip, bypasses the norm gate and is not evidence about the deployed pipeline. The second runs the shipped gate and aggregation code against five attacks at $K\in\{10,20\}$ and up to 40\% malicious clients, with eight seeds. Two results matter for the argument above. The gate rejects every sign-flipped update and no honest one, and is blind to the other four attacks. The median keeps the ranking quality of the global model at $K\ge 2f+1$ but not its operating point: at 40\% malicious clients the bound still holds ($K=10$, $f=4$ gives $2f+1=9$; $K=20$, $f=8$ gives $17$), yet the false-positive rate at $\tau=0.15$ reaches 0.24--0.67 under four attacks. The bound therefore guarantees bounded influence on each coordinate, not an unchanged decision threshold.

\subsection{Authentication Properties}
\label{sec:sec-auth}
An update passes, in order: enrolment lookup, constant-time MAC comparison, timestamp within $120$\,s, nonce unseen in $300$\,s, per-identity rate limit, dimension and finiteness check (HTTP~400), norm gate (HTTP~409), and finally aggregation. The first failure rejects the update. Unknown and malformed identities receive the same response, so the endpoint does not reveal which identifiers are enrolled. The replay nonce is consumed during verification, before the rate limiter runs, so an update that draws a 429 cannot be resent verbatim. HMAC gives authenticity and freshness; it does not hide the update from an observer (Section~\ref{sec:sec-limits}).

\subsection{Test Plan for the Control Plane}
\label{sec:sec-attacks}
\texttt{test\_auth\_attacks.py} contains ten cases, each built so that a correct control plane returns one specific status code, and run against a live aggregator. The harness rebuilds every message instead of resending it, so a retry past a 429 cannot trigger the replay guard and mask the case under test, and the one case meant to be accepted replays the current global weights, leaving the deployed model unchanged. The repository holds no archived output from a run. Table~\ref{tab:security:attacks} is therefore a test plan with expected verdicts, not a record of passes.

\begin{table}[t]
  \centering
  \caption{Planned control-plane test cases and the status code a correct aggregator should return (not an archived result).}
  \label{tab:security:attacks}
  \small
  \begin{tabular}{@{}p{3.0cm}@{\hspace{4pt}}c@{\hspace{4pt}}p{1.7cm}@{\hspace{4pt}}p{1.2cm}@{}}
    \toprule
    Case & Expect & Control & Threat \\
    \midrule
    Untagged update (old format); unknown self-tagged identity & 401 & enrolment, MAC & A2, A3 \\
    Forged tag; tampered weights with stale tag & 401 & MAC & A2 \\
    Replay of an accepted update; timestamp one hour old & 401 & nonce, clock & A2 \\
    Second update inside rate window & 429 & rate limit & A4 \\
    $50\times$ scaled weights & 409 & norm gate & A4 \\
    Wrong dimension; NaN weights & 400 & validation & A4 \\
    \bottomrule
  \end{tabular}
\end{table}

Even a clean run would show only that each control returns its designed verdict. It would not establish security, and it says nothing about plausible-norm, adversarially chosen weights from an enrolled node.

\subsection{Hardening Changes}
\label{sec:sec-audit}
A code review against the draft claims produced changes that bear on how earlier numbers are read. The original poisoning figure had been drawn from hard-coded values and was replaced by a measurement. The Byzantine argument had been void at $K=2$ and the update window was keyed by time rather than peer; both are fixed and $f$ is now reported. One model file replaces three inconsistent models. Handlers reject a missing or non-numeric \texttt{Content-Length} with HTTP~400, servers use \texttt{ThreadingHTTPServer} with a $5$\,s connection timeout, counters are updated under locks, and model snapshots are published by write, \texttt{fsync} and \texttt{os.replace}. Lowering the block threshold from $0.5$ to $0.15$ after data augmentation is an operating-point change, not a security fix; its false-positive cost is examined in Section~\ref{sec:eval-deployed}.

\subsection{Limitations}
\label{sec:sec-limits}
\paragraph{Byzantine tolerance at deployed scale.}
Two gateways give $f=0$ by~\eqref{eq:security:byz-bound}, and the poisoning results come from simulation. Authentication and the norm gate still apply at $K=2$, but neither stops a compromised gateway that uploads plausible-norm weights chosen adversarially.

\paragraph{Adaptive attacks.}
The simulation includes one aggregation-aware attack, the full-knowledge attack of Fang \textit{et al.}~\cite{fang2020}, and no adaptive timing, backdoors or collusion beyond a fixed strategy. The classifier is linear over hashed trigrams and simple syntactic features, and padding, entropy shaping or steering trigrams toward lightly weighted buckets can plausibly evade it. No input-space evasion was evaluated, so the poisoning results say nothing about evasion.

\paragraph{Authentication.}
Symmetric keys mean a breach of the aggregator exposes every peer secret, and the scheme offers no non-repudiation. Replay and rate-limit state live in aggregator memory and are lost on restart. Enrolment is manual, which bounds $f$ and also costs operator effort. The timestamp window assumes loosely synchronised clocks.

\paragraph{Confidentiality and privacy.}
Live updates and global models travel as plain-HTTP JSON on a trusted network segment, and the ML-KEM exchange is a standalone benchmark (Section~\ref{sec:sys-pqc}). An observer of the link can read the weights, and a recorded exchange could be read later by a quantum-capable adversary; A2's recording capability is not neutralised. Differential privacy is not implemented, so no $\epsilon$ is reported and gradient inversion was not measured. The parameters an update carries are a function of a site's training data, and what that reveals is unquantified.

\paragraph{Statistical scope.}
The attack tests assert designed status codes and were not run as a study. The bare-aggregator ablation is a single seeded run; the deployed-pipeline sweep has eight seeds and a simulated, omniscient adversary. What the analysis supports is narrow: authenticated, rate-limited, norm-checked updates from enrolled peers, and a median rule that preserves ranking quality under the tested attacks while $K\ge 2f+1$, at a cost in false positives that grows with the malicious fraction.
